% GENERATED FILE. Edit canonical lessons, then run npm run generate:books.
% canonical-source-hash: fnv1a64:5918ae43903ffbc9
% canonical-lessons: TE-C447-signalu, TE-C447-laisensu, TE-C447-padacari, TE-C447-kalibata, TE-C447-atikramincu

\chapter{Traffic Rules}
\label{ch:te-traffic-rules}
\hlchaptermodality{447}

\begin{chapteropening}
\textbf{By the end of this chapter:} \emph{I can say a traffic signal, a licence, a pedestrian, a footpath and to break a rule.}

Complete the last lesson of chapter 447: I can say a traffic signal, a licence, a pedestrian, a footpath and to break a rule.
\end{chapteropening}

\section[\texorpdfstring{\te{సిగ్నలు} (signalu)}{సిగ్నలు (signalu)}]{\te{సిగ్నలు} (signalu) — a traffic signal}

\label{lesson:TE-C447-signalu}



\begin{quote}
Before the new one: say the Telugu for to overturn, then the Telugu for to get dented.
\end{quote}

\subsection*{You'll want to know: \te{సిగ్నలు}}
\textbf{\te{సిగ్నలు}} — \emph{signalu} — \textquotedblleft{}a traffic signal\textquotedblright{}.

\begin{quote}
\te{ఎర్ర} \te{సిగ్నలు} \te{దగ్గర} \te{వాహనాలు} \te{ఆగాలి}.
\end{quote}

\begin{quote}
\emph{erra signalu daggara vāhanālu āgāli.}
\end{quote}

\begin{quote}
Vehicles must stop at a red signal.
\end{quote}

\te{ఎర్ర} \te{సిగ్నలు}, a red light; \te{పచ్చ} \te{సిగ్నలు}, a green light.

\begin{cousinweb}
\textbf{\te{సిగ్నలు}} (\emph{signalu}) is a loan from English \emph{signal}, with the Telugu final -\te{ఉ}.
\end{cousinweb}

\begin{grammarlens}[title={What you've built}]
One more word for getting around.
\end{grammarlens}

\subsection*{Your turn}
\begin{itemize}
\raggedright
  \item \emph{Say it:} \emph{signalu}
  \item \emph{Say it:} \emph{signalu}, once more
  \item \emph{Say it:} \emph{erra signalu daggara vāhanālu āgāli.} — Vehicles must stop at a red signal.
  \item \emph{Recall:} say the Telugu for to overturn, then the Telugu for to get dented, then say \emph{signalu} again
\end{itemize}

\begin{tcolorbox}[breakable,colback=teal!4,colframe=teal!35!black,title={Before you move on}]
What does \te{సిగ్నలు} mean? (\textquotedblleft{}A traffic signal\textquotedblright{}.) Say it once more.
\end{tcolorbox}
\section[\texorpdfstring{\te{లైసెన్సు} (laisensu)}{లైసెన్సు (laisensu)}]{\te{లైసెన్సు} (laisensu) — a licence}

\label{lesson:TE-C447-laisensu}



\begin{quote}
Before the new one: say the Telugu for to get dented, then the Telugu for a traffic signal.
\end{quote}

\subsection*{You'll want to know: \te{లైసెన్సు}}
\textbf{\te{లైసెన్సు}} — \emph{laisensu} — \textquotedblleft{}a licence\textquotedblright{}.

\begin{quote}
\te{లైసెన్సు} \te{లేకుండా} \te{బండి} \te{నడపకూడదు}.
\end{quote}

\begin{quote}
\emph{laisensu lēkuṇḍā baṇḍi naḍapakūḍadu.}
\end{quote}

\begin{quote}
You must not drive without a licence.
\end{quote}

A driving licence is \te{డ్రైవింగ్} \te{లైసెన్సు}; a shop needs a trade \te{లైసెన్సు} too.

\begin{cousinweb}
\textbf{\te{లైసెన్సు}} (\emph{laisensu}) is a loan from English \emph{licence}, with the Telugu final -\te{ఉ}.
\end{cousinweb}

\begin{grammarlens}[title={What you've built}]
One more word for getting around.
\end{grammarlens}

\subsection*{Your turn}
\begin{itemize}
\raggedright
  \item \emph{Say it:} \emph{laisensu}
  \item \emph{Say it:} \emph{laisensu}, once more
  \item \emph{Say it:} \emph{laisensu lēkuṇḍā baṇḍi naḍapakūḍadu.} — You must not drive without a licence.
  \item \emph{Recall:} say the Telugu for to get dented, then the Telugu for a traffic signal, then say \emph{laisensu} again
\end{itemize}

\begin{tcolorbox}[breakable,colback=teal!4,colframe=teal!35!black,title={Before you move on}]
What does \te{లైసెన్సు} mean? (\textquotedblleft{}A licence\textquotedblright{}.) Say it once more.
\end{tcolorbox}
\section[\texorpdfstring{\te{పాదచారి} (pādacāri)}{పాదచారి (pādacāri)}]{\te{పాదచారి} (pādacāri) — a pedestrian}

\label{lesson:TE-C447-padacari}



\begin{quote}
Before the new one: say the Telugu for a traffic signal, then the Telugu for a licence.
\end{quote}

\subsection*{You'll want to know: \te{పాదచారి}}
\textbf{\te{పాదచారి}} — \emph{pādacāri} — \textquotedblleft{}a pedestrian\textquotedblright{}.

\begin{quote}
\te{పాదచారులు} \te{రోడ్డు} \te{దాటుతున్నప్పుడు} \te{కార్లు} \te{ఆగాయి}.
\end{quote}

\begin{quote}
\emph{pādacārulu rōḍḍu dāṭutunnappuḍu kārlu āgāyi.}
\end{quote}

\begin{quote}
The cars stopped while the pedestrians were crossing the road.
\end{quote}

Plural: \te{పాదచారులు}, \emph{pādacārulu}.

\te{పాద}, foot, + -\te{చారి}, one who moves: one who goes on foot. Plural \te{పాదచారులు}.

\begin{cousinweb}
\textbf{\te{పాదచారి}} (\emph{pādacāri}) is borrowed from Sanskrit \emph{pādacārin}, \textquotedblleft{}going on foot\textquotedblright{} (Monier-Williams), from \emph{pāda}, \textquotedblleft{}foot\textquotedblright{}, and \emph{cārin}, \textquotedblleft{}moving\textquotedblright{}.
\end{cousinweb}

\begin{grammarlens}[title={What you've built}]
One more word for getting around.
\end{grammarlens}

\subsection*{Your turn}
\begin{itemize}
\raggedright
  \item \emph{Say it:} \emph{pādacāri}
  \item \emph{Say it:} \emph{pādacāri}, once more
  \item \emph{Say it:} \emph{pādacārulu rōḍḍu dāṭutunnappuḍu kārlu āgāyi.} — The cars stopped while the pedestrians were crossing the road.
  \item \emph{Recall:} say the Telugu for a traffic signal, then the Telugu for a licence, then say \emph{pādacāri} again
\end{itemize}

\begin{tcolorbox}[breakable,colback=teal!4,colframe=teal!35!black,title={Before you move on}]
What does \te{పాదచారి} mean? (\textquotedblleft{}A pedestrian\textquotedblright{}.) Say it once more.
\end{tcolorbox}
\section[\texorpdfstring{\te{కాలిబాట} (kālibāṭa)}{కాలిబాట (kālibāṭa)}]{\te{కాలిబాట} (kālibāṭa) — a footpath; a pavement}

\label{lesson:TE-C447-kalibata}



\begin{quote}
Before the new one: say the Telugu for a licence, then the Telugu for a pedestrian.
\end{quote}

\subsection*{You'll want to know: \te{కాలిబాట}}
\textbf{\te{కాలిబాట}} — \emph{kālibāṭa} — \textquotedblleft{}a footpath; a pavement\textquotedblright{}.

\begin{quote}
\te{పిల్లలు} \te{కాలిబాట} \te{మీదే} \te{నడవాలి}.
\end{quote}

\begin{quote}
\emph{pillalu kālibāṭa mīdē naḍavāli.}
\end{quote}

\begin{quote}
Children should walk only on the pavement.
\end{quote}

\te{కాలి}, of the foot, + \te{బాట}, a path. In a town it is the pavement beside the road.

\begin{cousinweb}
\textbf{\te{కాలిబాట}} (\emph{kālibāṭa}) is a Telugu compound: \te{కాలి}, the oblique of \te{కాలు}, \textquotedblleft{}foot, leg\textquotedblright{}, native Dravidian (DEDR 1479), and \te{బాట}, \textquotedblleft{}a path\textquotedblright{}, usually traced to Sanskrit \emph{vartman}, \textquotedblleft{}a track\textquotedblright{}, through Prakrit (CDIAL 11366).
\end{cousinweb}

\begin{grammarlens}[title={What you've built}]
One more word for getting around.
\end{grammarlens}

\subsection*{Your turn}
\begin{itemize}
\raggedright
  \item \emph{Say it:} \emph{kālibāṭa}
  \item \emph{Say it:} \emph{kālibāṭa}, once more
  \item \emph{Say it:} \emph{pillalu kālibāṭa mīdē naḍavāli.} — Children should walk only on the pavement.
  \item \emph{Recall:} say the Telugu for a licence, then the Telugu for a pedestrian, then say \emph{kālibāṭa} again
\end{itemize}

\begin{tcolorbox}[breakable,colback=teal!4,colframe=teal!35!black,title={Before you move on}]
What does \te{కాలిబాట} mean? (\textquotedblleft{}A footpath; a pavement\textquotedblright{}.) Say it once more.
\end{tcolorbox}
\section[\texorpdfstring{\te{అతిక్రమించు} (atikramiñcu)}{అతిక్రమించు (atikramiñcu)}]{\te{అతిక్రమించు} (atikramiñcu) — to break a rule; to violate}

\label{lesson:TE-C447-atikramincu}



\begin{quote}
Before the new one: say the Telugu for a pedestrian, then the Telugu for a footpath.
\end{quote}

\subsection*{You'll want to know: \te{అతిక్రమించు}}
\textbf{\te{అతిక్రమించు}} — \emph{atikramiñcu} — \textquotedblleft{}to break a rule; to violate\textquotedblright{}.

\begin{quote}
\te{అతను} \te{సిగ్నలు} \te{అతిక్రమించి} \te{జరిమానా} \te{కట్టాడు}.
\end{quote}

\begin{quote}
\emph{atanu signalu atikramiñci jarimānā kaṭṭāḍu.}
\end{quote}

\begin{quote}
He jumped the signal and paid a fine.
\end{quote}

Past (I …): \te{అతిక్రమించాను}, \emph{atikramiñcānu}.

A formal word. What is broken takes no ending: \te{నియమాలు} \te{అతిక్రమించు}, to break the rules.

\begin{cousinweb}
\textbf{\te{అతిక్రమించు}} (\emph{atikramiñcu}) is built from Sanskrit \emph{atikrama}, \textquotedblleft{}a going beyond, overstepping\textquotedblright{} (Monier-Williams), from \emph{ati}, \textquotedblleft{}beyond\textquotedblright{}, and \emph{kram}, \textquotedblleft{}to step\textquotedblright{}, with the Telugu verb ending -\te{ఇంచు}.
\end{cousinweb}

\begin{grammarlens}[title={What you've built}]
One more word for getting around.
\end{grammarlens}

\subsection*{Your turn}
\begin{itemize}
\raggedright
  \item \emph{Say it:} \emph{atikramiñcu}
  \item \emph{Say it:} \emph{atikramiñcu}, once more
  \item \emph{Say it:} \emph{atanu signalu atikramiñci jarimānā kaṭṭāḍu.} — He jumped the signal and paid a fine.
  \item \emph{Recall:} say the Telugu for a pedestrian, then the Telugu for a footpath, then say \emph{atikramiñcu} again
\end{itemize}

\begin{tcolorbox}[breakable,colback=teal!4,colframe=teal!35!black,title={Before you move on}]
What does \te{అతిక్రమించు} mean? (\textquotedblleft{}To break a rule; to violate\textquotedblright{}.) Say it once more.
\end{tcolorbox}
